\section{Application 1 Details}
\label{app:app1}

\subsection{Evaluation Details}
We evaluate Application 1 on six complementary benchmarks: AgentHarm, AgentSafetyBench, AgentDojo, AgentDyn, AgentSecurityBench, and BFCL. These benchmarks jointly measure harmful-request refusal, agentic tool-use safety, indirect prompt-injection robustness, task utility in interactive environments, and function-calling accuracy.

Unless otherwise specified, all models are evaluated with the same Qwen3.5-native chat and tool-call format used during training. We also disable model thinking during evaluation so that all models are compared under the same non-CoT inference setting. For benchmarks requiring automatic judging, we use the same LLM-as-judge model GLM-5.1~\citep{glm5team2026glm5} and judging prompt across all evaluated models.

\paragraph{AgentHarm.}
AgentHarm~\citep{agentharm2024} evaluates direct harmful-request refusal and benign-request preservation. We evaluate both the harmful and benign splits. For the harmful split, we report the harmfulness score (HS), where lower is better, and the refusal rate (RR), where higher is better. For the benign split, we report the benign score (BS) and the benign refusal rate to measure over-refusal. This benchmark is used to diagnose whether a model can reject explicitly harmful user requests without unnecessarily refusing safe agentic requests.

All models are evaluated with the same deterministic generation setting and the same AgentHarm evaluation scripts. The harmful and benign outputs are saved separately and summarized into the final model-level metrics. In the main table, AgentHarm therefore captures two complementary effects: improved safety on harmful requests and potential utility loss from over-refusal on benign requests.

\paragraph{AgentSafetyBench.}
AgentSafetyBench~\citep{zhang2024agentsafetybench} evaluates safety on risk-annotated agentic tool-use examples. The benchmark focuses on unsafe user intents, risky tool calls, unsafe tool arguments, and harmful task outcomes. For Qwen3.5 models, we use the native tool-call prompt described above: tool definitions are rendered into the prompt, model outputs are parsed as Qwen3.5-native tool calls, and thinking is disabled.

Each generated trajectory is judged by the AgentSafetyBench safety evaluator to determine whether the model response leads to an unsafe action or unsafe completion. We report the safe rate (SR), where higher is better; equivalently, this is one minus the unsafe rate. We do not include the legacy AgentSafetyBench helpfulness score in the main comparison, because the AgentSafetyBench evaluation set is primarily safety-oriented and does not provide a clean, benign split suitable for utility comparison.

\paragraph{AgentDojo.}
AgentDojo~\citep{agentdojo2024} evaluates robustness to indirect prompt injection in realistic tool-use environments. The benchmark contains benign user tasks with adversarial instructions embedded in untrusted external observations such as emails, webpages, and documents. We evaluate the model using the Qwen3.5-native parsed tool-call backend, so that tool calls are parsed from the same native XML-style format used in training.

For each trajectory, the evaluator records whether the original user task is completed and whether the injected malicious objective succeeds. We report benign utility (BU), utility under attack (UA), and attack success rate (ASR). BU and UA are higher-is-better metrics, while ASR is lower-is-better. AgentDojo is reported separately from AgentDyn because its task construction and environment dynamics differ from AgentDyn, and merging them would hide distinct failure modes.

\paragraph{AgentDyn.}
AgentDyn~\citep{li2026agentdyn} evaluates dynamic, stateful agentic environments with Shopping, GitHub, and DailyLife suites. Compared with AgentDojo, AgentDyn places more emphasis on longer interaction traces, state-changing tool calls, and task-dependent authorization. The same Qwen3.5-native parsed tool-call backend is used for AgentDyn as for AgentDojo.

We report benign utility (BU), utility under attack (UA), and attack success rate (ASR). AgentDyn is particularly useful for diagnosing whether a model can distinguish legitimate high-impact actions from injected or irrelevant side effects. Examples include payment actions, account updates, repository operations, file operations, and calendar scheduling. A good model should block malicious or irrelevant injected actions while preserving authorized actions that are necessary for completing the original user task.

\paragraph{AgentSecurityBench.}
AgentSecurityBench~\citep{zhang2025agent} evaluates full agent-security behavior under a broader set of attack and clean scenarios. We use the updated AgentSecurityBench protocol with the Qwen3.5-native tool-call format, disabled thinking, and no workflow-level case decomposition during runtime. Final case-level metrics are parsed only after the full evaluation run completes.

We do not evaluate the entire AgentSecurityBench benchmark. Instead, we use the sampled evaluation subset. This subset contains 3,035 evaluated rows across 17 evaluation shards. Among them, 2,716 rows are attack scenarios used to compute the primary AgentSecurityBench attack success rate (ASR), and the remaining 319 rows are smoke, clean, or clean-POT diagnostic scenarios used to monitor benign task behavior. The evaluated attack settings include direct prompt injection (DPI), indirect prompt injection (OPI), mixed attacks, persistent object threat (POT) backdoor attacks, memory-write attacks, and memory-read attacks. The clean or diagnostic settings include smoke DPI, POT clean, and clean no-attack tasks.

For attack scenarios, we report ASR, where lower is better. We also track original-task success, clean-task success, and refusal behavior as diagnostic metrics. In the main comparison table, we report AgentSecurityBench ASR as the primary security metric because it directly measures whether the model completes the adversarial objective on the evaluated attack subset.

\paragraph{BFCL.}
BFCL~\citep{patil2025bfcl} evaluates function-calling correctness rather than safety. We use the non-live function-calling categories currently configured on the cluster: \texttt{simple\_python}, \texttt{simple\_java}, \texttt{simple\_javascript}, \texttt{multiple}, \texttt{parallel}, \texttt{parallel\_multiple}, and \texttt{irrelevance}. The model is evaluated with the corresponding Qwen3.5-native function-calling setup, and the final score measures whether the model selects the correct function and generates correct arguments under BFCL's canonical schemas.

We report the overall BFCL accuracy as a capability metric. BFCL is shown separately from the safety benchmarks because it primarily measures schema adherence and function-call correctness, rather than refusal or robustness to adversarial instructions.

\paragraph{Metric Interpretation}
Across all safety benchmarks, lower harmfulness, unsafe rate, and ASR indicate stronger safety. Higher harmful-request refusal rate, benign score, task utility, and BFCL accuracy indicate better usefulness or tool-use capability. We treat AgentDojo and AgentDyn as separate benchmarks because they expose different robustness and utility failures. In the main analysis, we compare models along both safety and utility axes to identify whether a method improves refusal and attack robustness at the cost of over-refusal or degraded tool-use performance.

\subsection{Lightweight Environment Synthesis and Deployment}
\label{app:env}
A central challenge in agentic reinforcement learning lies in constructing interactive environments that yield reliable feedback. While realistic software environments provide the optimal feedback, fully replicating real-world environments is computationally expensive, difficult to scale, and often impractical for broadly deployable safety-alignment research. To address this, we design purpose-built and finite-state Python simulators that preserve the essential interaction dynamics required for RL as shown in Figure~\ref{fig:pipline_env}. By isolating only task-relevant resources, tool interfaces, and rule-based rewards, our approach trades strict real-world fidelity for practical deployability, computational efficiency, and scenario-specific reliability. Finally, we rigorously filter the generated tasks and environments for overall quality and code correctness, yielding a reliable utility dataset ready for the subsequent risk-injection phase.

Specifically, the utility synthesis process begins by sampling a subset of tools from a predefined pool. Based on these, we plan viable tasks by constructing tool-call graphs, which explicitly constrain and control the underlying task complexity. Guided by this graph-based plan, we define the essential environmental resources (e.g., mock files, emails) and trackable states, precisely mapping out which states and resources each tool is permitted to access or modify. Leveraging this well-defined scope, LLMs generate the underlying Python code that simulates the tools, while we concurrently formulate rule-based reward functions grounded in the expected resources and states. To guarantee robustness, every step is verified against predefined constraints, triggering automated repairs if requirements are unmet. Finally, post-generation, the complete task-environment-reward tuples undergo a rigorous filtering phase evaluating empirical task complexity, reward rationality, query naturalness, and code executability.

Built upon the core triad of \textit{environment}, \textit{task}, and \textit{reward}, safety environments are first synthesized from clean tasks to provide coherent initial states, tool affordances, and benign objectives for scalable agent rollouts. On top of these environments, attacked safety tasks are constructed by introducing adversarial risks while preserving executable tool-use dynamics, enabling agents to generate diverse trajectories under both benign and attacked conditions. The construction pipeline first identifies clean tasks with feasible tool affordances and risk-bearing execution paths, and then synthesizes paired clean and attacked scenarios together with structured rule-based feedback signals that distinguish benign task completion, harmful action execution, and safe refusal or confirmation-seeking behavior, thereby supporting downstream reward modeling and agentic safety reinforcement learning. The current framework mainly supports two complementary attack settings. 

For environment injection attacks, adversarial payloads are injected into contextual environment content such as documents, notes, or messages while the original user request remains benign, evaluating whether agents propagate corrupted contextual information into downstream tool actions. For malicious query attacks, the environment remains unchanged and the adversarial signal is instead introduced through malicious or partially malicious user requests that induce unsafe objectives such as unauthorized transfers or harmful state modifications. The above two scenarios capture adversarial environment manipulation and adversarial user intent, respectively. As illustrated in Figure~\ref{fig:env_scal}, our designed environments demonstrate remarkable scalability and resource efficiency. Even under a heavy workload, where the server is simultaneously loading 10000 environments, maintaining 1000 active environments, and executing 1,000 tool calls, the system maintains consistently stable latency, with the peak memory footprint strictly bounded below 2.5 GB.
